\documentclass[10pt,a4paper]{amsart}
\usepackage[margin=19mm]{geometry}
\usepackage{amsmath,amssymb,mathtools}
\usepackage[scr=rsfs,cal=euler]{mathalfa}
\usepackage{xfrac}
\usepackage[unicode,hidelinks,bookmarks=false]{hyperref}
\newcommand{\ie}{i.e.\ }
\newcommand{\cf}{cf.\ }
\newcommand{\resp}{respectively\ }
\newcommand{\Subsection}{\subsection}
\providecommand{\nobreakdash}{\nobreak\hbox{-}}
\setlength{\emergencystretch}{2em}
\multlinegap0pt
\usepackage{amssymb}
\newtheorem{theorem}{Theorem}
\newtheorem{lemma}{Lemma}
\newcommand{\ud}{\mathrm{d}}
\title{Abrupt decorrelation for linear stochastic differential equations}
\author{Sergio I. L\'opez, Juan C. Pardo, Leandro P. R. Pimentel}
\date{Source: arXiv:2509.16828v2 (CC BY 4.0)}
\begin{document}
\maketitle

\noindent\emph{Source and licence.} Abrupt decorrelation for linear stochastic differential equations. Version arXiv:2509.16828v2 (CC BY 4.0), distributed by arXiv under the Creative Commons Attribution 4.0 International licence, http://creativecommons.org/licenses/by/4.0/, as recorded in the arXiv licence metadata for this exact version and shown on the abstract page https://arxiv.org/abs/2509.16828v2. Full licence notice: \texttt{licenses/arxiv-2509.16828-CC-BY-4.0.txt}.\par\smallskip

\noindent\textit{Editorial scope.} Counted unit: the article's lemma Jordan (source0.tex lines 716--724). The article's complete original proof of it is delivered with it (source0.tex lines 725--769, one \texttt{proof} environment of 45 lines). No mathematics is written, repaired or re-derived: the statement and the proof are the article's own lines, in the article's own notation, and no retained line is substituted or re-flowed.  Retained uncounted as the source-local context the proof needs: lines 107--109; lines 358--401.  Nothing was omitted from any retained span, so the package carries no omission record.  No retained line contains a citation, so the wrapper carries no bibliography and no bibliography entry.  No retained line contains a figure environment, an image-inclusion command or drawing code, so no image is delivered and the package declares no asset dependency.  Editorial formatting only, in addition: an AMS \texttt{amsart} preamble that restates the article's own declarations and macros used by the retained lines, namely the environments \texttt{theorem}, \texttt{lemma} on separate counters, together with the standard packages those lines need.  The retained environments print in this document as Theorem 1, Lemma 1 in order of appearance, whereas the article prints the same items with the numbering of its own full document.  The complete native source of the article as distributed in the arXiv e-print for this exact version is bundled unchanged as \texttt{sources/185415/source0.tex}.
\begin{equation}\label{SDE}
\ud X^{(\epsilon)}_t =QX^{(\epsilon)}_t \ud t +\epsilon \sigma \ud B_t\qquad \textrm{ for any }\quad t>0,
\end{equation}

\begin{theorem}\label{theorem2}  
Consider the family of OU processes $(X^{(\epsilon)}; \epsilon>0)$, defined in \eqref{SDE} and assume that  $Q$ is Hurwitz stable and the controllability rank condition is fulfilled. Then, abrupt decorrelation occurs with respect to both the  Kullback-Liebler divergence and in the total variation distance at the time scale
$$  t_\epsilon:=\frac{1}{|\vartheta|}\Big(|\ln \epsilon|+\left(m-1\right)\ln|\ln \epsilon|\Big).$$  Moreover,
\begin{itemize}
	\item[i)] when we are in the complex case of the leading eigenvalues $\lambda^*$'s, the family $(X^{(\epsilon)}; \epsilon>0)$ exhibits
	window decorrelation with asymptotically constant window sizes $w_\epsilon\to \frac{w}{|\vartheta|}>0$ in both Kullback-Leibler divergence and total variation distance;
	
	\item[ii)] in the case where there is a unique real eigenvalue $\lambda^*$, the family $\{X^{(\epsilon)}; \epsilon>0\}$ exhibits decorrelation profile phenomenon with  asymptotically constant window sizes $w_\epsilon\to \frac{w}{|\vartheta|}>0$ under; 
	\begin{enumerate}
		\item[a)] The Kullback-Leibler divergence with profile function
		\[
		\begin{split}
			G_{KL}(r) &= \lim_{\epsilon\to0}D_{KL}( \mathcal{L}^{(\epsilon)}_{t_\epsilon+r\omega_{\epsilon}} \Big|  \mathcal{I}^{(\epsilon)}_{t_\epsilon+r\omega_{\epsilon}} )  \\
			& =   -\frac{1}{2}\ln\left(\det\left( \mathbf{I}_d- \Gamma\Sigma_0\Gamma^T\left(\Gamma\Sigma_0\Gamma^T+ e^{2rw}\varsigma_{\infty} \right)^{-1}\right) \right)\,,
		\end{split}
		\]
		where 
		\[
		\Gamma=\frac{|\vartheta|^{1-m}}{(m-1)!} P_1\mathbf{I}^{\ell}_mP_1^{-1}, 
		\] 
		the matrix $P_1$ is obtained by rearranging the block order of $P$ and 
		\[
		\mathbf{I}^{\ell}_m=\begin{pmatrix}
			E_{1m}^{(m)}& & & & &\\ 
			&\ddots & & & &\\
			& &E_{1m}^{(m)} & & &\\
			& & & 0& & \\
			& & & &\ddots & \\
			& & & & & 0
		\end{pmatrix},
		\]
		with $\ell\geq 1$ denoting the number of times that $E_{1m}^{(m)}$ appears along the diagonal of $\mathbf{I}^{\ell}_m$;
		\item[b)] The total variation distance with profile function
		$$G_{TV}(r):=\lim_{\epsilon\to0}||\mathcal{L}^{(\epsilon)}_{t_\epsilon+r\omega_{\epsilon}}-\mathcal{I}^{(\epsilon)}_{t_\epsilon+r\omega_{\epsilon}}||_{TV} = ||\mathcal{N}(0, M(r) )-\mathcal{N}(0,\mathbf{I}_{2d} )||_{TV}, $$ 
		where $M(r)$ is a block matrix given by 
		\[
		M(r)=\begin{pmatrix}
			\mathbf{I}_d &  \Sigma_0^{\frac{1}{2}} \Gamma^T  ( \Gamma \Sigma_0 \Gamma^T + e^{2r w} \varsigma_{\infty} )^{-\frac{1}{2}} 
			\\ ( \Gamma \Sigma_0 \Gamma^T + e^{2rw} \varsigma_{\infty} )^{-\frac{1}{2}} \Gamma \Sigma_0^{\frac{1}{2}} & \mathbf{I}_d  \end{pmatrix}\,.
		\]
	\end{enumerate} 
\end{itemize}

\end{theorem}

\begin{lemma}\label{Jordan} Let $Q\in\mathbb{R}^{d\times d}$ be a Hurwitz-stable matrix.  Denote by $\vartheta$  the real part of any critical eigenvalue, and let $t_\epsilon$  and  $\omega_\epsilon$ be defined as in the statement of Theorem \ref{theorem2}. \begin{itemize}
		\item[i)] If there is a unique real critical eigenvalue, then the following limit holds
		\[
		\lim_{\epsilon \downarrow 0} \frac{e^{(t_{\epsilon}+r\omega_\epsilon) Q}}{\epsilon}=e^{-rw}\frac{|\vartheta|^{1-m}}{(m-1)!} P_1\mathbf{I}^{\ell}_mP_1^{-1}=e^{-rw}\Gamma\,,
		\]
		where $P_1$, $\mathbf{I}^{\ell}_m$ and $\Gamma$ are defined as in the statement of Theorem \ref{theorem2}.
		\item[ii)] In the complex case, the above  limit does not converge in general. However, there exist a subsequence $(\epsilon_n)_{n\ge 1}$ with $\epsilon_n\downarrow 0$ such that the limit  exists along this subsequence.
	\end{itemize}
\end{lemma}

\begin{proof} From the Jordan decomposition of $Q$, we have  $e^{tQ}=PJP^{-1}$, where
	$$J=\begin{pmatrix}
		e^{tB_1}& & \\ 
		&\ddots &  \\
		& & e^{tB_p} 
	\end{pmatrix},$$
	is the Jordan canonical form of $e^{tQ}$. Observe that all the others entries in $J$ are zero and each $e^{tB_j}$ is a square matrix of one of the following forms
	\begin{equation}\label{eq:Jordan}
		(i)\,\, e^{tB_j}=\begin{pmatrix}
			e^{\lambda_j t}&te^{\lambda_j t}&\dots&\frac{t^{m_j-1}}{(m_j-1)!}e^{\lambda_j t}\\ 
			&e^{\lambda_j t}&\dots&\frac{t^{m_j-2}}{(m_j-2)!}e^{\lambda_j t}\\
			& &\ddots & \vdots \\
			& & & e^{\lambda_j t } 
		\end{pmatrix}\,\mbox{ or } (ii)\,\, e^{tB_j}=\begin{pmatrix}e^{\Lambda_j t}&te^{\Lambda_j t}&\dots&\frac{t^{m_j-1}}{(m_j-1)!}e^{\Lambda_j t}\\ 
			&e^{\Lambda_j t}&\dots&\frac{t^{m_j-2}}{(m_j-2)!}e^{\Lambda_j t}\\
			& &\ddots & \vdots \\
			& & & e^{\Lambda_j t }\end{pmatrix}
	\end{equation}
	with 
	$$e^{\Lambda_j t}=e^{\alpha_jt} \begin{pmatrix}\cos\left(\beta_j t\right)&-\sin\left(\beta_j t\right)\\
		\sin\left(\beta_j t\right)& \cos\left(\beta_j t\right)\end{pmatrix}\,.$$
	This representation follows from the Jordan decomposition, where in (i) the eigenvalue $\lambda_j\in\mathbb{R}$ and $e^{tB_j}$ is a $m_j\times m_j$ matrix,  while in (ii) the eigenvalue $\lambda_j=\alpha_j+i\beta_j\in\mathbb{C}$ has a non zero imaginary part and $e^{tB_j}$ is a $2m_j\times 2m_j$ matrix. Since we are assuming that 
	\[
	\vartheta=\max \{ \Re(\lambda_i)\,:\,j=1,\dots,k \} <0\,,
	\] 
	it happens that $e^{tQ}\to 0$ as $t\to\infty$. In order to capture the block $e^{tB_j}$
	that decays at a slower rate, we note that there may exist multiple blocks associated with 
	$\vartheta$.  Let
	$$m:=\max\{m_j\,:\,\Re(\lambda_i)=\vartheta\}\geq 1\,, $$ 
	that is, $m$ is the largest among the $m_j$ associated with $\vartheta$, so that $t^{m-1}e^{\vartheta t}$  corresponds to the slowest decaying term appearing in $e^{tQ}$. Let  $t_\epsilon$  and  $\omega_\epsilon$ be as defined  in the statement of Theorem \ref{theorem2}.
	Then 
	$$\left(t_\epsilon + r w_\epsilon\right)^{m-1} e^{\vartheta(t_\epsilon + r w_\epsilon)}=\left(t_\epsilon + r w_\epsilon\right)^{m-1}\frac{\epsilon}{\left(\log\frac{1}{\epsilon}\right)^{m-1}}e^{-rw+\mathcal{O}(\epsilon)}\,,$$
	and since 
	$$\lim_{\epsilon\downarrow 0} \frac{\left(t_\epsilon + r w_\epsilon\right)^{m-1}}{\left(\log\frac{1}{\epsilon}\right)^{m-1}}=|\vartheta|^{1-m}\,,$$
	it follows that 
	$$\lim_{\epsilon\downarrow 0}\frac{\left(t_\epsilon + r w_\epsilon\right)^{m-1} e^{ \vartheta(t_\epsilon + r w_\epsilon)}}{\epsilon} =|\vartheta|^{1-m}e^{-rw}\,.$$
	If we have a block $e^{-tB_j}$ associated to $\vartheta$ and $m$ that has the real form (i), that is $\lambda_j=\vartheta$ and $m_j=m$, then
	$$\lim_{\epsilon\downarrow 0}\frac{ e^{ (t_\epsilon+rw_\epsilon) B_j}}{\epsilon}= e^{-rw}\frac{|\vartheta|^{1-m}}{(m-1)!} E^{(m)}_{1m}\,,$$
	where we recall that $E^{(m)}_{1m}$ is the $m\times m$ matrix with a 1 in the $(1,m)$-th entry and zeros elsewhere.
	
	If we have a block $e^{tB_j}$ associated to $\vartheta$ and $m$ that has the complex form (ii), that is $\alpha_j=\vartheta$ and $m_j=m$, then the respective rotation matrix (with $\sin\left(\beta_jt\right)$ and $\cos\left(\beta_j t\right)$) will stay bounded but will not converge, although one can see that we do have limits along subsequences. 
	
	Finally, all  blocks not associated with $\vartheta$ and $m$  converge to zero. This implies  that  the limit of $\epsilon^{-1}e^{ \left(t_\epsilon + r w_\epsilon\right) Q}$, as $\epsilon\downarrow 0$, exists  if and only if all  blocks associated with $\vartheta$ and $m$ are of the form (i),  assuming that there are $\ell\ge 1$ such blocks.  By appropriately rearranging, we can assume that all block associated with $\vartheta$ and $m$ appear first along the diagonal. Under this arrangement, the limit can be expressed as
	$$e^{-rw}\frac{|\vartheta|^{1-m}}{(m-1)!} P_1\mathbf{I}^{\ell}_mP_1^{-1}.$$   
	This completes the proof. \end{proof}

\end{document}
